% Outline: Background (1.5 page cap, >=2 figures; defines the paper's vocabulary)
\section{Background}

\subsection{Serving, Scheduling, and Head-of-Line Blocking}
Modern LLM engines admit requests continuously and batch them at
token granularity; continuous batching~\cite{orca2022} and vLLM's paged
KV-cache~\cite{vllm2023} are the substrate for this work. The scheduler decides which waiting request joins
the running batch next. First-come-first-served is the default, and its
weakness is classical: a long-running request admitted early delays every
short request behind it. A 2026 survey of latency in LLM agent systems
independently lists FIFO-induced head-of-line blocking among live system
bottlenecks for agentic traffic~\cite{park2026minimizing}, while placing
learned scheduling in its future-work section --- the gap this paper works in.

\subsection{Learned Output-Length Prediction}
Shortest-job-style ordering requires knowing output lengths in
advance. Since generation length is stochastic, a body of work predicts it:
S\textsuperscript{3} classifies responses into length buckets; the
learning-to-rank line~\cite{fuLTR2024} showed that \emph{relative order},
not absolute length, is what scheduling needs, training a small proxy with a
ranking loss; PARS~\cite{tao2026ranking} refined this with pairwise margin
training over a BERT encoder. All of these read the prompt (or scalar request
statistics); none consume tool-schema text.

\subsection{System Substrate: Gateway and Decision Service}
Our serving stack (Fig.~\ref{fig:arch}) fronts the engine with an
LLM gateway (VeloxMesh) that owns routing, admission, and a scheduler hook
with a documented default timeout: scoring must return within 15\,ms
(operator-configurable), after which
the gateway \emph{fails open} to arrival order. The gateway consults a
separate decision service over HTTP; that service
hosts the Ranker and applies the Reliability Gate on the gateway's behalf.
This fail-open contract is load-bearing for everything that follows: an
unreliable or slow prediction never blocks the request path.

\begin{figure*}[t]
  \centering
  % ONE grid, ONE module.
  % Columns: 0--52 | 62--114 | 124--176 -- three 52 mm columns and a single
  % 10 mm gutter. Rows: +18..-18 and -26..-62 -- two 36 mm containers and a
  % single 8 mm row gap, so every container on the page shares one top edge
  % and one bottom edge. Inside every container: 3 mm of side padding and 2 mm
  % of top and bottom padding, a header strip of one fixed 6 mm height pinned
  % top-left, and a 46x24 mm content region cut into a 9 mm body lane at
  % yc+3.5, a 6 mm connector lane, and a 9 mm body lane at yc-11.5. Every body
  % box is 46x9 mm, and every connector starts and ends exactly on the border
  % of the box it names (style: edge gap = 0).
  \begin{tikzpicture}[x=1mm,y=1mm]
    \begin{scope}[on background layer]
      \node[panel, minimum width=52mm, minimum height=36mm] (client) at (26,0) {};
      \node[panel, minimum width=52mm, minimum height=36mm] (gw)     at (88,0) {};
      \node[panel, minimum width=52mm, minimum height=36mm] (engine) at (150,0) {};
      \node[panel accent, minimum width=52mm, minimum height=36mm] (dec) at (88,-44) {};
      \node[panel, minimum width=52mm, minimum height=36mm] (verd) at (150,-44) {};
    \end{scope}

    % ---- header strips: one language, one height, one baseline per row ------
    \node[tier title, anchor=north west] at (3,16)    {Agent client};
    \node[tier title, anchor=north west] at (65,16)   {Gateway};
    \node[tier title, anchor=north west] at (127,16)  {vLLM engine};
    \node[tier title, anchor=north west] at (65,-28)  {Decision service};
    \node[tier title, anchor=north west] at (127,-28) {Verdict};

    % ---- row 1, col 1: the one captured turn every number in the paper uses --
    \node[environment component, inner sep=0pt, minimum width=46mm,
          minimum height=9mm] (cli) at (26,3.5) {OpenCode};
    \node[environment component, align=center, text width=42mm, inner sep=0pt,
          minimum width=46mm, minimum height=9mm] (turn) at (26,-11.5)
      {Captured turn\\{\sub{170 tools, 147\,kB}}};
    \draw[data flow, edge gap] (26,-1) -- (26,-7);

    % ---- row 1, col 2: gateway (one flush strip, then the scoring hook) ------
    \node[environment component, inner sep=0pt, minimum width=46mm,
          minimum height=9mm] (strip) at (88,3.5) {};
    \draw[draw=black!45, line width=.7pt] (80.333,-1) -- (80.333,8);
    \draw[draw=black!45, line width=.7pt] (95.667,-1) -- (95.667,8);
    \node at (72.667,3.5) {Admit};
    \node at (88,3.5)     {Cache};
    \node at (103.333,3.5){Route};
    \node[environment component, align=center, text width=42mm, inner sep=0pt,
          minimum width=46mm, minimum height=9mm] (hook) at (88,-11.5)
      {Scheduler hook\\{\sub{verdict due in 15\,ms}}};
    \draw[data flow, edge gap] (88,-1) -- (88,-7);

    % ---- row 1, col 3: the queue lives INSIDE vLLM, not beside it -----------
    % (verified: the gateway attaches the verdict to the forwarded request,
    %  VeloxMesh internal/ltr/decision.go:218; the policy reads it back from
    %  the request's extra_args inside the engine, vllm_scheduler.py:87)
    \node[proposed component, align=center, text width=42mm, inner sep=0pt,
          minimum width=46mm, minimum height=9mm] (sched) at (150,3.5)
      {Scheduler\\{\sub{reorders trusted slots}}};
    \node[trusted slot, minimum width=10mm, minimum height=9mm] (q1) at (132,-11.5) {};
    \node[pinned slot,  minimum width=10mm, minimum height=9mm] (q2) at (144,-11.5) {};
    \node[trusted slot, minimum width=10mm, minimum height=9mm] (q3) at (156,-11.5) {};
    \node[pinned slot,  minimum width=10mm, minimum height=9mm] (q4) at (168,-11.5) {};
    % Both ends land on a slot's top border; the swap needs no chip because the
    % scheduler box above it already says what the swap is.
    \draw[data flow, edge gap]
      (156,-7) .. controls (150,-2) and (138,-2) .. (132,-7);

    % ---- row 2, col 2: decision service, one flush three-lane stack ---------
    \node[proposed component, inner sep=0pt, minimum width=46mm,
          minimum height=24mm] (stack) at (88,-48) {};
    \draw[draw=sysaccent!85!black, line width=.9pt] (65,-45) -- (111,-45);
    \draw[draw=sysaccent!85!black, line width=.9pt] (65,-51) -- (111,-51);
    \node at (88,-40.5) {Stratum classifier};
    \node at (88,-48) {Reliability Gate};
    \node at (88,-55.5) {Schema-aware Ranker};

    % ---- row 2, col 3: the verdict, as two boxes an arrow can land on -------
    \node[environment component, align=center, text width=42mm, inner sep=0pt,
          minimum width=46mm, minimum height=9mm] (fo) at (150,-40.5)
      {Arrival order\\{\sub{abstain, or 15\,ms passes}}};
    \node[proposed component, align=center, text width=42mm, inner sep=0pt,
          minimum width=46mm, minimum height=9mm] (key) at (150,-55.5)
      {Reorder key\\{\sub{trusted: S3 / S4}}};
    \draw[degraded flow, edge gap, lane turn]
      (111,-48) -- (119,-48) -- (119,-40.5) -- (127,-40.5);
    \draw[data flow, edge gap] (111,-55.5) -- (127,-55.5);

    % ---- the two inter-row connectors: ask down, verdict up ----------------
    % Each chip sits beside its own arrow, and the arrow is offset by half the
    % chip so that arrow+chip centre on the column axis (x=88 and x=150).
    \draw[data flow, edge gap] (75.75,-18) -- (75.75,-26);
    \node[metric badge, anchor=west] at (77.25,-22) {score, 15\,ms};
    \draw[data flow, edge gap] (141.75,-26) -- (141.75,-18);
    \node[metric badge, anchor=west] at (143.25,-22) {verdict};

    % ---- edges on the serving path -----------------------------------------
    \draw[data flow, edge gap, lane turn] (49,-11.5) -- (57,-11.5) -- (57,3.5) -- (65,3.5);
    \draw[data flow, edge gap] (111,3.5) -- (127,3.5);

    % ---- legend: a compact band, deliberately smaller than a subsystem ------
    \node[tier band, minimum width=52mm, minimum height=16mm, inner sep=0pt]
      at (26,-44) {};
    \node[proposed component, inner sep=0pt, minimum width=7mm,
          minimum height=4.4mm] at (8.5,-39) {};
    \node[fact label, fill=none, anchor=west, text=black!75] at (13.5,-39)
      {added by this work};
    \node[environment component, inner sep=0pt, minimum width=7mm,
          minimum height=4.4mm] at (8.5,-44) {};
    \node[fact label, fill=none, anchor=west, text=black!75] at (13.5,-44)
      {unchanged substrate};
    \draw[degraded flow] (5,-49) -- (12,-49);
    \node[fact label, fill=none, anchor=west, text=black!75] at (13.5,-49)
      {fail-open path};
  \end{tikzpicture}
  \caption{The gateway scores at admission inside a 15\,ms budget and
  forwards the verdict; the engine's scheduler reorders trusted slots among
  themselves and leaves abstained ones where they arrived. A timeout or an
  abstention yields arrival order.}
  \label{fig:arch}
\end{figure*}

\begin{figure*}[t]\centering
  % Source: runs/coldstart/ (999 test requests, 3 seeds) and the gate's
  % validation fit; every number is read from those artifacts at build time.
  \includegraphics[width=\linewidth]{strata.pdf}
  \caption{The four cold-start strata, and which of them the paper is
  allowed to speak for. S1 and S2 hold too few test requests to support a
  claim, so their Ranking Tau is withheld and the gate routes them to
  Fallback.}
  \label{fig:strata}
\end{figure*}

\subsection{Vocabulary}
Throughout, the \emph{Ranker} is the fine-tuned BERT model scoring a
request's expected output length; the \emph{Reliability Gate} is the runtime
decision of whether that score may influence queue order; \emph{Fallback}
means serving in arrival order because the gate withheld trust.
\emph{Ranking Tau} is Kendall's $\tau_b$~\cite{kendall1938} between predicted and true
output-length order on held-out data. \emph{Cold-Start Transfer} measures
Ranking Tau on test strata whose tool sets were unseen in training, within
the same workload family: S1 (seen combination), S2 (new combination of seen
tools), S3 (partially new tools), S4 (all-new tools); Fig.~\ref{fig:strata}
gives each stratum's size and what the gate does with it. It is distinct from
\emph{Cross-Workload Transfer}, which changes the workload family entirely.
